% Generated draft: requires mathematical and source review.
\documentclass[11pt]{article}
\usepackage[margin=25mm]{geometry}
\usepackage{fontspec}
\setmainfont{Latin Modern Roman}
\usepackage{amsmath,amssymb,mathtools,unicode-math}
\setmathfont{Latin Modern Math}
\usepackage{enumitem,needspace,xurl,hyperref}
\setlength{\parindent}{0pt}
\setlength{\parskip}{5pt}
\setlength{\emergencystretch}{3em}
\newcommand{\N}{\mathbb{N}}
\newcommand{\Z}{\mathbb{Z}}
\newcommand{\Q}{\mathbb{Q}}
\newcommand{\R}{\mathbb{R}}
\newcommand{\C}{\mathbb{C}}
\newcommand{\F}{\mathbb{F}}
\newcommand{\A}{\mathbb{A}}
\newcommand{\PP}{\mathbb{P}}
\newcommand{\E}{\mathbb{E}}
\newcommand{\eps}{\varepsilon}
\newcommand{\dd}{\,\mathrm{d}}
\begin{document}
\section*{UnsolvedMath ID: 30001021}
\textbf{Source problem code:} catalogue entry
\subsection{Definitions and mathematical statement}
Let $f:X\to Y$ be a surjective morphism of smooth projective varieties with connected fibers, an algebraic fiber space as in the source. Set \[\omega_{X/Y}=\omega_X\otimes f^*\omega_Y^{-1},\qquad E_m=f_*\omega_{X/Y}^{\otimes m}=f_*\mathcal O_X(mK_{X/Y}).\] Let $U\subseteq Y$ denote the complement of the source's discriminant locus; in the usual characteristic-zero interpretation it is the open set over which $f$ is smooth. Global generation outside that locus means that the evaluation map using sections on all of $Y$ is surjective there: \[H^0(Y,E_m)\otimes\mathcal O_U\longrightarrow E_m|_U\quad\text{is surjective}.\] The question is whether \[\exists m_0\geq1\ \forall m\in\Z_{\geq m_0}:\qquad E_m\text{ is globally generated on }U\text{ by its global sections on }Y.\] This retains every sufficiently large integer $m$, not only sufficiently divisible ones. It also retains the untwisted sheaf: no ample line-bundle tensor factor is supplied. Generation by sections defined merely over $U$, weak positivity, generic generation after twists or covers, and this evaluation-map surjectivity are different properties and are not identified. A zero sheaf satisfies the evaluation condition vacuously; the source does not require nonzero pluricanonical spaces. The ground field and discriminant convention are not expressly specified.\par The supplied 2026-08-21 triage labels the record partially solved and describes an almost global generation theorem and a resulting Kodaira-dimension inequality, while saying Iitaka subadditivity remains open in general with further special cases. It does not give that theorem's hypotheses, twists, or divisibility conditions. This theorem wording therefore does not verify the stronger-looking untwisted, every-large-$m$ question as written, nor supply missing assumptions; both the displayed question and the literature qualification are retained.
\subsection{Short English statement}
Must every sufficiently high relative pluricanonical pushforward be generated off the discriminant by sections on the whole base? The statement has no ample twist and says all sufficiently large integers, so it must be distinguished from weak positivity or sufficiently divisible generation results.
\subsection{Sources}
\begin{itemize}\item Local UnsolvedMath registry, permanent ID 30001021, source problem catalogue entry, \emph{Global Generation of Relative Pluricanonical Pushforwards}, proposed year 2008; no individual proposer supplied.\item Supplied citation: \emph{Komplexe Analysis} (2009), Oberwolfach Reports, DOI 10.4171/owr/2008/38, pp.\ 2165--2218: \url{https://doi.org/10.4171/owr/2008/38}. Supplied literature label: original report.\item Supplied literature label: Chen and Han, generalized nonvanishing and Iitaka conjecture, \url{https://arxiv.org/abs/2312.04015}.\item Literature assessment date: 2026-08-21. The status and progress claims are supplied dated curation, not independently verified here.\end{itemize}
\end{document}
